% GENERATED FILE. Edit canonical lessons, then run npm run generate:books.
% canonical-source-hash: fnv1a64:3df827380e0f86be
% canonical-lessons: HI-C152-khun, HI-C152-dakghar, HI-C152-pustakalay, HI-C152-mandir, HI-C152-masjid

\chapter{Blood, Post office, Library, Temple, Mosque}
\label{ch:hi-blood-post-office-library-temple-mosque}
\hlchaptermodality{152}

\begin{chapteropening}
\textbf{By the end of this chapter:} \emph{I can say blood, a post office, a library, a temple and a mosque.}

Complete the last lesson of chapter 152: I can say blood, a post office, a library, a temple and a mosque.
\end{chapteropening}

\section[\texorpdfstring{\hi{ख़ून} (khūn)}{ख़ून (khūn)}]{\hi{ख़ून} (khūn) — blood}

\label{lesson:HI-C152-khun}



\begin{quote}
Before the new one: say the Hindi for the chest, then the Hindi for the waist.
\end{quote}

\subsection*{You'll want to know: \hi{ख़ून}}
\textbf{\hi{ख़ून}} — \emph{khūn} — \textquotedblleft{}blood\textquotedblright{}.

\begin{cousinweb}
\textbf{\hi{ख़ून}} (\emph{khūn}) is a loan, usually traced to Persian \emph{khūn}, \textquotedblleft{}blood\textquotedblright{}.
\end{cousinweb}

\begin{grammarlens}[title={What you've built}]
One more word for this chapter.
\end{grammarlens}

\subsection*{Your turn}
\begin{itemize}
\raggedright
  \item \emph{Say it:} \emph{khūn}
  \item \emph{Say it:} \emph{khūn}, once more
  \item \emph{Say it:} say \emph{khūn}
  \item \emph{Recall:} say the Hindi for the chest, then the Hindi for the waist, then say \emph{khūn} again
\end{itemize}

\begin{tcolorbox}[breakable,colback=teal!4,colframe=teal!35!black,title={Before you move on}]
What does \hi{ख़ून} mean? (\textquotedblleft{}Blood\textquotedblright{}.) Say it once more.
\end{tcolorbox}
\section[\texorpdfstring{\hi{डाकघर} (ḍākghar)}{डाकघर (ḍākghar)}]{\hi{डाकघर} (ḍākghar) — a post office}

\label{lesson:HI-C152-dakghar}



\begin{quote}
Before the new one: say the Hindi for the waist, then the Hindi for blood.
\end{quote}

\subsection*{You'll want to know: \hi{डाकघर}}
\textbf{\hi{डाकघर}} — \emph{ḍākghar} — \textquotedblleft{}a post office\textquotedblright{}.

\begin{cousinweb}
\textbf{\hi{डाकघर}} (\emph{ḍākghar}) is \textquotedblleft{}post-house\textquotedblright{}: \emph{ḍāk}, \textquotedblleft{}the post\textquotedblright{}, from a reconstructed \emph{*ḍakka}, \textquotedblleft{}relay, post\textquotedblright{}, which CDIAL asks may belong with \emph{*ḍakk}, \textquotedblleft{}shout\textquotedblright{} (CDIAL 5519), with \emph{ghar}, \textquotedblleft{}house\textquotedblright{}, from \emph{ghara} (CDIAL 4428).
\end{cousinweb}

\begin{grammarlens}[title={What you've built}]
One more word for this chapter.
\end{grammarlens}

\subsection*{Your turn}
\begin{itemize}
\raggedright
  \item \emph{Say it:} \emph{ḍākghar}
  \item \emph{Say it:} \emph{ḍākghar}, once more
  \item \emph{Say it:} say \emph{ḍākghar}
  \item \emph{Recall:} say the Hindi for the waist, then the Hindi for blood, then say \emph{ḍākghar} again
\end{itemize}

\begin{tcolorbox}[breakable,colback=teal!4,colframe=teal!35!black,title={Before you move on}]
What does \hi{डाकघर} mean? (\textquotedblleft{}A post office\textquotedblright{}.) Say it once more.
\end{tcolorbox}
\section[\texorpdfstring{\hi{पुस्तकालय} (pustakālay)}{पुस्तकालय (pustakālay)}]{\hi{पुस्तकालय} (pustakālay) — a library}

\label{lesson:HI-C152-pustakalay}



\begin{quote}
Before the new one: say the Hindi for blood, then the Hindi for a post office.

Say the four mind verbs and what each first meant: \textbf{\hi{लिखना}}, to scratch; \textbf{\hi{पढ़ना}}, to recite; \textbf{\hi{समझना}}, to wake up to; \textbf{\hi{सोचना}}, to burn.
\end{quote}

\subsection*{You'll want to know: \hi{पुस्तकालय}}
\textbf{\hi{पुस्तकालय}} — \emph{pustakālay} — \textquotedblleft{}a library\textquotedblright{}.

\begin{cousinweb}
\textbf{\hi{पुस्तकालय}} (\emph{pustakālay}) is a tatsama compound, taken over whole from Sanskrit: \emph{pustaka}, \textquotedblleft{}book\textquotedblright{}, and \emph{ālaya}, \textquotedblleft{}abode\textquotedblright{}.
\end{cousinweb}

\begin{grammarlens}[title={What you've built}]
One more word for this chapter.
\end{grammarlens}

\subsection*{Your turn}
\begin{itemize}
\raggedright
  \item \emph{Say it:} \emph{pustakālay}
  \item \emph{Say it:} \emph{pustakālay}, once more
  \item \emph{Say it:} say \emph{pustakālay}
  \item \emph{Recall:} say the Hindi for blood, then the Hindi for a post office, then say \emph{pustakālay} again
\end{itemize}

\begin{tcolorbox}[breakable,colback=teal!4,colframe=teal!35!black,title={Before you move on}]
What does \hi{पुस्तकालय} mean? (\textquotedblleft{}A library\textquotedblright{}.) Say it once more.
\end{tcolorbox}
\section[\texorpdfstring{\hi{मंदिर} (mandir)}{मंदिर (mandir)}]{\hi{मंदिर} (mandir) — a temple}

\label{lesson:HI-C152-mandir}



\begin{quote}
Before the new one: say the Hindi for a post office, then the Hindi for a library.

Name the week by its sky: \textbf{\hi{सोमवार}}, the Moon's day, through \textbf{\hi{शनिवार}}, Saturn's, and \textbf{\hi{रविवार}}, the Sun's. It is the same seven-body week Rome kept.
\end{quote}

\subsection*{You'll want to know: \hi{मंदिर}}
\textbf{\hi{मंदिर}} — \emph{mandir} — \textquotedblleft{}a temple\textquotedblright{}.

\begin{cousinweb}
\textbf{\hi{मंदिर}} (\emph{mandir}) is Sanskrit \emph{mandira}, \textquotedblleft{}house, palace, temple\textquotedblright{}, taken over whole as a learned loan, a tatsama.
\end{cousinweb}

\begin{grammarlens}[title={What you've built}]
One more word for this chapter.
\end{grammarlens}

\subsection*{Your turn}
\begin{itemize}
\raggedright
  \item \emph{Say it:} \emph{mandir}
  \item \emph{Say it:} \emph{mandir}, once more
  \item \emph{Say it:} say \emph{mandir}
  \item \emph{Recall:} say the Hindi for a post office, then the Hindi for a library, then say \emph{mandir} again
\end{itemize}

\begin{tcolorbox}[breakable,colback=teal!4,colframe=teal!35!black,title={Before you move on}]
What does \hi{मंदिर} mean? (\textquotedblleft{}A temple\textquotedblright{}.) Say it once more.
\end{tcolorbox}
\section[\texorpdfstring{\hi{मस्जिद} (masjid)}{मस्जिद (masjid)}]{\hi{मस्जिद} (masjid) — a mosque}

\label{lesson:HI-C152-masjid}



\begin{quote}
Before the new one: say the Hindi for a library, then the Hindi for a temple.
\end{quote}

\subsection*{You'll want to know: \hi{मस्जिद}}
\textbf{\hi{मस्जिद}} — \emph{masjid} — \textquotedblleft{}a mosque\textquotedblright{}.

\begin{cousinweb}
\textbf{\hi{मस्जिद}} (\emph{masjid}) is a loan from Arabic \emph{masjid}, \textquotedblleft{}place of prostration, mosque\textquotedblright{}, which reached Hindi through Persian. Its Arabic root is \emph{s-j-d} (Lane).
\end{cousinweb}

\begin{grammarlens}[title={What you've built}]
One more word for this chapter.
\end{grammarlens}

\subsection*{Your turn}
\begin{itemize}
\raggedright
  \item \emph{Say it:} \emph{masjid}
  \item \emph{Say it:} \emph{masjid}, once more
  \item \emph{Say it:} say \emph{masjid}
  \item \emph{Recall:} say the Hindi for a library, then the Hindi for a temple, then say \emph{masjid} again
\end{itemize}

\begin{tcolorbox}[breakable,colback=teal!4,colframe=teal!35!black,title={Before you move on}]
What does \hi{मस्जिद} mean? (\textquotedblleft{}A mosque\textquotedblright{}.) Say it once more.
\end{tcolorbox}
